\subsection{A caveat on the standard Q-Learning: it depends on the order of the actions}
\label{sec:order}
The standard Q-Learning of our batteries behaved in ways that looked like learning but were not. On the standard benchmark it lost to random choice by up to $1.19$ points in cells where the static choice gains most (Section~\ref{sec:fanjul}); on the published families of unrelated machines it gained up to $65\%$ of the validated gain (Section~\ref{sec:families}); and at equal CPU time it won ``by choosing cheaper actions'' (Section~\ref{sec:rq2}). A diagnostic of the standard benchmark (Entry~258, registered before it was run) showed what the learner actually does: it puts $55$--$73\%$ of its choices on the action listed first (random choice: $1/A$, between $3\%$ and $20\%$), while UCB1 and the policy gradient stay at the share of random choice ($3$--$20\%$). Table~\ref{tab:order_a} gives the shares.

% Auto-generated by analysis/make_table_journal.py
\begin{table}[t]
\caption{Order sensitivity of the standard Q-Learning on the standard benchmark ($140$ test instances, one per cell of type, $m$ and $n$; two seeds; gains over random choice in points, positive = better). ql: standard; ql\_v0: the same learner with the class's validated action moved to index $0$; val: the validated action. Index-0 share: fraction of the learner's choices that fall on the action listed first (random choice: $1/A$); ql\_v0 share: fraction on the validated action.}
\label{tab:order_a}
\centering
{\scriptsize\setlength{\tabcolsep}{5pt}
\begin{tabular}{@{}lrrrrrr@{}}
\toprule
Space & $1/A$ & ql index-0 share & ql & ql\_v0 & val & ql\_v0 share of val \\
\midrule
B5 & 0.20 & 0.73 & $-$0.36 & 0.73 & 0.90 & 0.74 \\
B15 & 0.07 & 0.71 & $-$0.02 & 1.30 & 1.62 & 0.70 \\
B30 & 0.03 & 0.55 & $-$0.54 & 2.02 & 2.35 & 0.68 \\
\bottomrule
\end{tabular}}
\end{table}


The first causal test, reversing the order of the actions, was badly designed: it made the learner lock on the last action of the list, another extreme and poor action, so it did not help (three of the seven predictions of that stage held). The second test moves the validated action of each class to index $0$ without changing anything else. The standard Q-Learning then gains $0.73$, $1.30$ and $2.02$ points over random choice in the three spaces, $81\%$, $81\%$ and $86\%$ of the validated gain (all five registered predictions of this stage held), and it plays the validated action in $68$--$74\%$ of its choices. Standing on the first index, it recovers most of the best fixed action; standing on any other, it recovers an average one. The learner's results depend on where the good action was listed.

\paragraph{The order-invariant learner.} We repeated the cells where the standard learner seemed to gain or lose with a version that permutes the actions at random in each run (\textsc{ql-perm}; Table~\ref{tab:order_b}). On the standard benchmark it gains over random choice beyond the margin in none of the $21$ cells (mean $-0.30$ points, against $-0.25$ for the standard learner); on the three published families it loses to random choice in all six cells (from $-0.05$ to $-0.66$ points; the standard learner's gains of $0.67$ and $0.83$ points with correlated machines were $53\%$ and $65\%$ of the validated gain, those of the permuted learner are negative), and the validated action beats it in all six; on the synthetic $P||C_{\max}$ battery it gains in none of the $24$ cells and the validated action beats it in $23$. Four of the six registered predictions on the standard benchmark and the families held, and all four on $P||C_{\max}$.

% Auto-generated by analysis/make_table_journal.py
\begin{table}[t]
\caption{The order-invariant Q-Learning (ql\_perm: a random permutation of the actions in each run) against the standard one and the validated action: mean gain over random choice across the cells of each battery (points, positive = better) and number of cells with a gain beyond the equivalence margin. P$||C_{\max}$: $12$ classes $\times$ spaces S5 and S15.}
\label{tab:order_b}
\centering
{\scriptsize\setlength{\tabcolsep}{4pt}
\begin{tabular}{@{}lrrrrr@{}}
\toprule
Battery & cells & ql & ql\_perm & val & cells with ql\_perm gain \\
\midrule
Standard benchmark (Fanjul-Peyro, Ruiz) & 21 & $-$0.25 & $-$0.30 & 1.60 & 0 \\
Published $R||C_{\max}$ families & 6 & 0.28 & $-$0.33 & 1.28 & 0 \\
Synthetic $P||C_{\max}$ & 24 & 0.00 & $-$0.06 & 0.09 & 0 \\
\bottomrule
\end{tabular}}
\end{table}


\paragraph{The cost effect.} In the $P||C_{\max}$ battery the actions are listed by destruction size, so the first action is the cheapest. The standard learner evaluates $0.44$--$0.46$ times as many solutions as random choice per run; the order-invariant learner evaluates $0.99$--$1.01$ times as many. The statement that at equal CPU time Q-Learning wins only by choosing cheaper actions is therefore true of the standard implementation, and its cause is the lock-in on the first index, not a policy learned from the reward.

\paragraph{What this changes.} It does not change the main conclusion; it makes it stronger: an online Q-Learning does not learn a useful policy in these budgets, and the largest gains attributed to it came from a good action standing first in the list. It changes three things in how the result must be read. First, ``Q-Learning gains in four of six cells of the published families'' is an artefact of the order of the actions and should not be read as evidence that online learning pays there. Second, the protocol needs a rule it did not have: report the learner's order sensitivity (here by permuting the actions) or break ties at random. Third, in the earlier batteries the standard learner performs like a single fixed action placed arbitrarily in the list; the optimistic Q-Learning, UCB1 and the other bandits, which do not repeat the first index, are the more reliable representatives of online learning and are the ones the claims rest on.
